\section{Related Work}

\subsection{Traditional Light Field Depth Estimation Methods}

Early light field depth estimation methods primarily extract geometric cues from the physical properties of light rays, with Epipolar Plane Image (EPI) slope analysis being a dominant approach. Wanner et al. formulate depth estimation as an orientation estimation problem in the EPI space. By computing the structure tensor of EPI lines, this method extracts local disparity values based on the dominant gradient direction. A continuous global optimization framework incorporating Markov Random Fields (MRF) is subsequently employed to enforce spatial smoothness and resolve depth ambiguities. This approach demonstrates high accuracy in regions with distinct textures and clear linear EPI structures, establishing a foundational baseline for EPI-based geometry extraction.

While EPI slope analysis performs effectively in textured regions, it degrades significantly in textureless areas where gradient information is sparse. To address this limitation, Tao et al. [2] integrate defocus cues with correspondence cues to construct a more robust depth estimation pipeline. The correspondence cue relies on photo-consistency across different viewpoints, similar to traditional stereo matching, whereas the defocus cue measures the sharpness of refocused images at varying depth planes. By combining these two complementary physical signals through a probabilistic framework, the method yields reliable depth maps even on weakly textured surfaces, expanding the applicability of light field depth estimation beyond purely texture-dependent EPI slope calculations.

Advancing the multi-view geometric perspective, Jeon et al. tackle the challenge of low angular resolution in commercial light field cameras, which restricts the baseline for conventional stereo matching. They introduce a sub-pixel phase shift technique to align multi-view images accurately, enabling the computation of high-precision correspondence costs. By incorporating adaptive support weights and handling occlusions explicitly, this method enhances depth estimation in regions with fine details and narrow structures. The integration of phase-shifted multi-view stereo with light field data effectively mitigates the quantization errors inherent in discrete angular sampling.

Despite their theoretical elegance, these traditional methods share significant limitations rooted in their reliance on handcrafted features and the strict Lambertian reflectance assumption. First, they assume that surface radiance remains invariant across viewpoints, mathematically forming linear EPI structures as $I(x, u) \approx f(x + u \cdot d)$. When scenes contain non-Lambertian materials, such as specular reflections or translucent surfaces, viewpoint-dependent radiance shifts destroy the linear EPI topology, causing the structure tensor and photo-consistency metrics to fail entirely [3]. Second, handcrafted features like gradients and structure tensors blur in complex texture regions and lack the capacity to explicitly decouple material properties from geometric depth. Consequently, these methods cannot distinguish whether an EPI distortion originates from a depth discontinuity or a surface reflection. Finally, existing approaches process mixed-material scenes uniformly, lacking a systematic root-cause analysis and physical boundary definition for EPI assumption failures. Unlike these traditional frameworks that suffer from severe performance degradation in non-Lambertian regions, our work introduces a three-layer angular signal decomposition model to decouple geometric features from material properties. By explicitly defining the physical boundaries of EPI validity and employing a dual-mask routing architecture, the proposed method adaptively handles mixed-material scenes, preventing the ineffective computational waste associated with applying pure EPI frameworks to non-Lambertian surfaces.

\subsection{Deep Learning-based General Light Field Depth Estimation Methods}

Building upon the physical cue extraction discussed in the previous section, deep learning-based methods have shifted the paradigm from handcrafted geometric descriptors to data-driven feature learning. Early convolutional neural network (CNN) architectures primarily focus on extracting disparity cues directly from Epipolar Plane Images (EPIs). For instance, Shin et al. [1] propose a multi-scale CNN that projects the 4D light field into 2D EPI slices to regress depth values. By learning the mapping between EPI slope patterns and disparity, this approach achieves high accuracy in textured Lambertian regions. However, the single-branch architecture implicitly assumes a uniform reflectance model, making the learned EPI features highly sensitive to viewpoint-dependent radiance variations caused by non-Lambertian surfaces.

To compensate for the limited receptive field of local EPI patches, subsequent research integrates sub-aperture images (SAIs) and macropixel features to capture broader spatial-angular correlations. Heber et al. [2] introduce a framework that jointly processes EPIs and macropixels, fusing local geometric slopes with global spatial context. More recently, Transformer-based architectures have been adopted to model long-range dependencies within the light field. The Light Field Transformer (LFT) [4] utilizes epipolar-shifted window attention to aggregate features across distant viewpoints, effectively reducing depth ambiguity in textureless areas. Despite these architectural advancements, the feature fusion remains implicit. The network treats all angular signals uniformly, failing to distinguish between the linear structures of Lambertian regions and the distorted patterns of specular or translucent materials.

Recognizing the vulnerability of standard EPI features to complex textures and reflectance, several methods incorporate attention mechanisms and adaptive loss functions to enhance robustness. For example, hierarchical attention networks [2] assign dynamic weights to different angular views, attempting to suppress inconsistent viewpoints that violate the photo-consistency assumption. Similarly, confidence map estimation approaches [2] down-weight unreliable disparity predictions in occluded or highly reflective areas. While these strategies yield moderate improvements in controlled environments, they rely on statistical regularization rather than physical modeling. The network is forced to learn a compromised representation that averages the characteristics of both Lambertian and non-Lambertian signals within a single feature space.

Despite the continuous evolution from basic CNNs to complex Transformers, these general light field depth estimation methods share significant limitations when applied to real-world mixed scenes. First, single-branch architectures typically assume the scene consists of pure Lambertian surfaces or simple mixtures, leading to significant performance degradation in complex urban environments where Lambertian and non-Lambertian materials coexist. Second, these methods lack explicit perception and feature routing mechanisms for non-Lambertian regions, causing the network to produce severe depth artifacts in specular highlights and fail to achieve material-adaptive processing. Third, implicit feature learning cannot explain the physical root causes of EPI failure. Empirical evidence indicates that merely increasing network complexity, introducing attention modules, or adjusting loss functions cannot overcome the physical limits imposed by the violation of the EPI assumption. These inherent bottlenecks necessitate a paradigm shift from implicit single-branch learning to explicit physical modeling. This motivates our unified dual-mask physical model, which decouples material properties via geometric features to enable a dual-mask routing architecture for mixed scenes, while systematically defining the physical boundaries of EPI failure to guide future research beyond pure EPI frameworks.

\subsection{Deep Learning-based Non-Lambertian and Complex Scene Light Field Depth Estimation Methods}

To mitigate the degradation of EPI structures in complex and mixed-material scenes, recent studies have explored multi-branch and dual-branch architectures that disentangle spatial and angular features. For instance, Wang et al. [2] propose a disentangled spatial-angular network that employs separate branches to extract local spatial textures and global angular correlations, significantly improving depth accuracy in occluded and textured regions. Building upon this, dual-branch frameworks specifically targeting material variations, such as the specular-aware network by Chen et al. [5], utilize an auxiliary branch to predict material masks and route features accordingly. While these methods demonstrate improved robustness in mixed scenes, their mask learning mechanisms are predominantly end-to-end black boxes. They lack a strict physical decomposition of the angular signal, making it difficult to extract discriminative geometric features for accurate material classification under low angular resolution settings (e.g., $9 \times 9$).

Complementing feature disentanglement, other approaches incorporate physical priors and auxiliary cues, such as defocus or specular highlights, to compensate for the failure of the Lambertian assumption. Huang et al. introduce a physics-guided framework that jointly estimates depth and specular masks, leveraging defocus cues to regularize depth predictions in highly reflective regions. By explicitly modeling the specular component, this method yields considerable improvements on synthetic datasets with prominent reflections. However, the integration of such auxiliary cues substantially increases the computational burden. More importantly, these methods fail to systematically define the physical boundaries where the EPI assumption breaks down. When applied to real-world scenarios where non-Lambertian training samples are inherently scarce, the reliance on auxiliary physical priors often leads to severe overfitting or performance bottlenecks, as the model cannot reliably distinguish between genuine physical reflections and complex Lambertian textures.

Departing from explicit EPI and SAI representations, Neural Radiance Fields (NeRF) and implicit neural representations have been adapted for light field depth estimation to inherently handle view-dependent non-Lambertian effects. Methods like LF-NeRF and Plenoxels [6] parameterize the light field using multi-layer perceptrons or sparse voxel grids, effectively modeling complex reflectance and specularities without relying on linear EPI structures. Although these implicit methods achieve high-fidelity view synthesis and robust depth estimation in purely non-Lambertian environments, they exhibit notable limitations in mixed urban or complex scenes. Specifically, they lack effective sampling strategies to address the multi-domain data imbalance inherent in mixed scenes. As a result, the models exhibit restricted generalization capabilities on minority domains (i.e., non-Lambertian regions) and struggle to achieve cross-domain exposure balance, often over-smoothing depth boundaries between Lambertian and non-Lambertian surfaces.

Despite these advancements, existing methods share substantial limitations in handling mixed non-Lambertian environments. First, current dual-branch approaches lack a rigorous physical decomposition of angular signals (e.g., decoupling into DC, linear, and residual components). Their mask learning mechanisms are predominantly end-to-end black boxes, making it difficult to extract discriminative geometric features for material classification under low angular resolutions (e.g., $9 \times 9$). Second, the integration of auxiliary cues, such as defocus, often increases computational burden. More importantly, these methods fail to systematically define the physical boundaries where the EPI assumption breaks down, leading to severe overfitting or performance bottlenecks when non-Lambertian training samples are scarce. Finally, existing frameworks lack effective sampling strategies to address multi-domain data imbalance in mixed scenes, which restricts model generalization on minority domains (non-Lambertian regions) and hinders cross-domain exposure balance. To address these gaps, this work proposes a unified dual-mask physical model. By introducing a three-layer angular signal decomposition, we decouple geometric and material features at the physical level. Our dual-mask and geometric routing architecture explicitly handles mixed materials, while our systematic analysis of the physical boundaries of EPI failure provides essential negative findings, preventing ineffective computational expenditure and guiding future research directions.

While existing multi-branch architectures improve robustness in mixed-material scenes, their predominantly black-box mask learning mechanisms lack a strict physical decomposition of angular signals and fail to uncover the fundamental root causes of Epipolar Plane Image (EPI) degradation on non-Lambertian surfaces. Inspired by these observations, we propose a unified dual-mask physical modeling approach that addresses the aforementioned limitations by systematically defining the physical boundaries of EPI validity, decoupling material and geometric features via a three-layer angular signal decomposition, and employing an explicit geometry-routed dual-branch network for robust depth estimation in complex scenarios. Building on these theoretical foundations, we now formalize these principles into a comprehensive computational framework, which is detailed in the following section.